% META: {"id":"1SPE-GEOREP-CO-020","chapitre":"1SPE-GEOMETRIE-REPEREE","type_objet":"corrige","exercice_ref":"1SPE-GEOREP-EX-020","capacites_codes":["C2"],"sources_inspiration":[],"status":"approved","origin":"generated","status_history":["generated","audited","accepted","approved"],"release_acceptance":"RELEASE_OWNER_BATCH_ACCEPTANCE","acceptance_closure_digest":"sha256:9b3ccf9a81c5520fbb7e03b7b2d3e7bf2057a02834b6d908bf3be5f49f3b553f","acceptance_receipt_digest":"sha256:4fad86f0282e0fa4e0419cca1ad6379ae7af5a280c04a4a90880f90a1c044242"}
% BEGIN-VERIFY
% from sympy import *
% x, y = symbols('x y')
% # Montrer que le triangle est rectangle
% # A(1,2), B(4,6), C(8,3)
% # AB = (3,4), AC = (7,1)
% # AB.AC = 21 + 4 = 25 != 0
% # AB = (3,4), BC = (4,-3)
% # AB.BC = 12 - 12 = 0 => angle droit en B
% assert 3*4 + 4*(-3) == 0
% END-VERIFY
\begin{corrige}{1SPE-GEOREP-EX-020}

\commentaireMarge{On calcule les produits scalaires de vecteurs issus de chaque sommet.}

$\vec{BA}\begin{pmatrix} -3 \\ -4 \end{pmatrix}$, $\vec{BC}\begin{pmatrix} 4 \\ -3 \end{pmatrix}$.

$\vec{BA} \cdot \vec{BC} = (-3)(4) + (-4)(-3) = -12 + 12 = 0$.

Donc $\vec{BA} \perp \vec{BC}$, le triangle $ABC$ est \textbf{rectangle en $B$}.

\end{corrige}
